%==============================================================================
\section{Introduction}
%==============================================================================

When an activist investor discovers underperformance at a portfolio company, she faces a consequential choice: sell her stake and move on, engage quietly behind the scenes, or accumulate shares aggressively enough to cross a regulatory disclosure threshold. Each path carries different risks and rewards---and each sends different signals to the market. How she resolves this dilemma shapes not only her own returns but whether a takeover bid emerges and what minority shareholders ultimately receive.

This tension is the classic exit-versus-voice problem of \citet{Hirschman1970}, transplanted into financial markets. On one side, liquidity facilitates governance by letting activists accumulate stakes cheaply and reducing free-riding \citep{Maug1998}. On the other, liquidity makes selling painless, tempting blockholders to exit rather than bear the costs of engagement \citep{Coffee1991,Bhide1993}. The tension deepens once we recognize that prices are not passive scorekeepers: when takeover decisions depend on stock prices, trading has real effects through feedback channels \citep{BondEdmansGoldstein2012,EdmansGoldsteinJiang2012}, so the blockholder's choice between exit and voice reverberates beyond her portfolio into the real economy.

No existing framework combines these three forces---endogenous governance choice, order-flow inference, and takeover feedback---in a single model. Yet they are tightly intertwined in practice. Activist campaigns frequently culminate in mergers and acquisitions, and takeover expectations help rationalize the returns activists earn \citep{BravJiangPartnoyThomas2008,GreenwoodSchor2009}. The free-rider problem that \citet{GrossmanHart1980} identified in tender offers makes takeover premia a natural welfare object for minority shareholders. Understanding how liquidity shapes these premia requires a model that takes all three forces seriously.

I develop a model in which a blockholder observes a private signal about firm value and chooses among four actions: Exit (sell her stake), Passive Hold (retain without engaging), Quiet Voice (engage below the disclosure threshold), or Public Voice (buy additional shares, engage, and trigger mandatory disclosure). A competitive market maker sets prices from discrete order flow and the disclosure flag, while a potential bidder conditions entry on these market signals. Engagement raises standalone value and strengthens resistance to acquisition, increasing the premium a bidder must offer. The key modeling choice is that disclosure is stake-triggered: activism is directly observable only when the blockholder crosses a regulatory threshold, creating two distinct information regimes within a single equilibrium.

The model predicts that liquidity has a non-monotone effect on expected minority takeover gains. At low liquidity, noise trading is scarce, so the blockholder's orders are highly informative; the market maker largely infers her action, deterring quiet engagement and suppressing bids. As liquidity rises, noise camouflages the activist's trades, encouraging voice and improving bid incidence. But at high liquidity, exit becomes so cheap that engagement incentives collapse entirely. The net effect is hump-shaped: minority shareholders benefit most at intermediate liquidity, where voice is both incentive-compatible and partially concealed.

Stake-triggered disclosure attenuates this inference channel. When the blockholder crosses the regulatory threshold, her engagement becomes directly observable rather than inferred from noisy order flow. I show that lowering the disclosure threshold shifts probability mass from Quiet Voice to Public Voice, compressing the activism premium into the observable regime and making it less sensitive to liquidity. This result speaks directly to policy debates: the U.S.\ Schedule 13D threshold at 5\% and the U.K.\ FCA threshold at 3\% create meaningfully different information environments, and the model provides a framework for evaluating how such differences affect takeover outcomes and minority shareholder welfare.

These predictions connect to well-documented empirical patterns. \citet{BravJiangPartnoyThomas2008} show that hedge fund activism often targets firms that subsequently receive acquisition offers. \citet{GreenwoodSchor2009} find that a substantial fraction of activist returns is attributable to takeover outcomes. Cross-country variation in disclosure stringency---from the relatively permissive U.S.\ regime to the stricter U.K.\ rules---provides a natural laboratory for testing the model's comparative statics on how disclosure thresholds shape the interaction between liquidity, activism, and corporate control.

Section~\ref{sec:lit} reviews the related literature. Section~\ref{sec:model} presents the model. Section~\ref{sec:equilibrium} characterizes equilibrium and derives the main results. Section~\ref{sec:numerical} provides numerical illustrations. Section~\ref{sec:comp-statics} presents comparative statics. Section~\ref{sec:extensions} discusses extensions, and Section~\ref{sec:conclusion} concludes.
